\documentclass[11pt]{article}
\usepackage[margin=1in]{geometry}
\usepackage[utf8]{inputenc}
\usepackage[T1]{fontenc}
\usepackage{lmodern}
\usepackage{amsmath,amssymb,amsthm}
\usepackage{booktabs}
\usepackage[title]{appendix}
\usepackage[colorlinks=true,linkcolor=blue,citecolor=blue,urlcolor=blue,pdfauthor={P. Graham},pdftitle={A four-dimensional 4-reptile simplex}]{hyperref}
\newtheorem{theorem}{Theorem}
\newtheorem{corollary}[theorem]{Corollary}
\newtheorem{lemma}[theorem]{Lemma}
\theoremstyle{remark}
\newtheorem{remark}[theorem]{Remark}
\setlength{\emergencystretch}{3em}
\title{A four-dimensional 4-reptile simplex}
\author{P. Graham\\ \small\url{https://orcid.org/0009-0007-8493-1505}}
\date{November 11, 2026}
\begin{document}
\maketitle
\begin{abstract}
A $d$-dimensional simplex is a $k$-reptile if it can be tiled by $k$
mutually congruent simplices similar to it. Matoušek and Safernová
conjectured that for $d\ge3$ a $k$-reptile $d$-simplex exists only when
$k=m^d$ for an integer $m$, and proved this for $d=3$. Kynčl and Patáková
proved that in dimension four $k$ must be a square, noted that every
$k$-reptile simplex known in dimension $d\ge3$ has $k=m^d$, and asked
whether a four-dimensional $k$-reptile simplex exists with $k$ a square
but not a fourth power. We show that the simplex with vertices $(0,0,0,0)$, $(1,0,0,0)$,
$(\tfrac12,\tfrac12,0,0)$, $(\tfrac23,\tfrac13,\tfrac13,0)$ and
$(\tfrac34,\tfrac14,\tfrac14,\tfrac14)$, the fundamental alcove of the
affine Weyl group of type $F_4$, is a $4$-reptile: it is tiled by four
congruent copies of itself scaled by $1/\sqrt2$, two of them mirror
images. This answers the question of Kynčl and Patáková with $k=4$ and
disproves the conjecture of Matoušek and Safernová in dimension four. The
subdivision itself is due to Felikson and Tumarkin, who found in 2004 that
the affine reflection group of type $F_4$ has a reflection subgroup of
index $4$ whose fundamental simplex is similar to the alcove. What we add
is the observation that this settles the reptile question, a short direct
proof of the tiling, and exact computer checks of it.

This work has not yet been reviewed by experts.
The proofs, the computer checks and the text were produced by AI models
(from OpenAI and Anthropic), working under the author's direction.
\end{abstract}
\medskip
\noindent\textit{Keywords:} reptile simplex, self-similar tiling,
alcove, affine Weyl group, $F_4$.

\noindent\textit{MSC 2020:} 52C22, 05B45, 20F55.

\section{Introduction}

\subsection{Main result}
A tiling of a closed set $X\subset\mathbb{R}^d$ is a decomposition of $X$
into finitely many closed sets with nonempty, pairwise disjoint
interiors. Following Kynčl and Patáková \cite{KP}, a $d$-dimensional
simplex $S$ is a \emph{$k$-reptile} if it can be tiled by $k$ simplices
that are all mutually congruent and similar to $S$; as they emphasize,
``congruence includes mirror symmetries'' \cite[p.~1]{KP}. Throughout, let
\begin{gather*}
S=\operatorname{conv}\{v_0,v_1,v_2,v_3,v_4\},\qquad\text{where}\\
v_0=(0,0,0,0),\quad v_1=(1,0,0,0),\quad v_2=(\tfrac12,\tfrac12,0,0),\quad
v_3=(\tfrac23,\tfrac13,\tfrac13,0),\quad v_4=(\tfrac34,\tfrac14,\tfrac14,\tfrac14).
\end{gather*}

\begin{theorem}\label{thm:main}
The simplex $S\subset\mathbb{R}^4$ is a $4$-reptile. More precisely, $S$
is the union of four simplices $S_0,\dots,S_3$ with pairwise disjoint
interiors, each the image of $S$ under a similarity of ratio $1/\sqrt2$.
Two of the four similarities preserve orientation and two reverse it.
\end{theorem}

Since $4=2^2$ and $4$ is not a fourth power, Theorem~\ref{thm:main}
answers the question of Kynčl and Patáková stated below, with $k=4$, and
gives a counterexample, in dimension $d=4$, to the conjecture of Matoušek
and Safernová.

The simplex $S$ is the fundamental alcove of the affine Weyl group of type
$F_4$ (Section~\ref{sec:origin}). A linear map $\varphi$ that is $\sqrt2$
times an isometry carries the hyperplane arrangement of this group into
itself, so $\varphi(S)$ is a union of alcoves, $4$ of them by volume.
Section~\ref{sec:proof} proves Theorem~\ref{thm:main} directly, by
writing $\varphi(S)$ and its four pieces as solution sets of linear
inequalities; the proof needs no Lie theory and no computer.
Section~\ref{sec:verification} describes exact computer checks of the
same tiling.

\subsection{The question}
Hertel \cite[pp.~650--652]{He} proved that the simplices he calls
H-simplices, and more generally his H-simplotopes, are $m^d$-replicating
(tiled by $m^d$ congruent copies of themselves scaled down) for every
integer $m\ge2$. His Conjecture~2 \cite[p.~652]{He} asserts that, for
$d>2$, an $f$-replicating $d$-simplex must have $f=m^d$ for an integer
$m\ge2$ and must be an H-simplex. Theorem~\ref{thm:main} also refutes this
conjecture in dimension four, since $4$ is not a fourth power.

Matoušek and Safernová \cite[Conjecture~1.2]{MS} conjectured: ``If $d\ge3$
and a $d$-dimensional $k$-reptile simplex exists, then $k=m^d$ for some
natural number $m\ge2$.'' They proved it for $d=3$; Kynčl and Patáková
gave a simpler proof \cite[Theorem~1.2]{KP} and proved that a
four-dimensional $k$-reptile simplex can exist only when $k$ is the square
of a positive integer \cite[Theorem~1.3]{KP}. In
their abstract they note that ``the only $k$-reptile simplices that are
known for $d\geqslant3$, have $k=m^d$''. They then ask \cite[Problem~1.4;
Problem~2 in the arXiv version]{KP} whether there is a four-dimensional
$k$-reptile simplex with $k$ the square of a non-square integer, that is,
with $k$ a square but not a fourth power. Theorem~\ref{thm:main} answers
this with $k=4$, the square of $2$. It is consistent with \cite[Theorem~1.3]{KP}, since
$k=4$ is a square, and it leaves the case $d=3$ untouched.

\subsection{Credit}
The tiling is not new as geometry. Felikson and Tumarkin
\cite[Lemma~10]{FT} consider Coxeter simplices $F\ne P$ in $\mathbb{E}^n$
with the same Coxeter diagram, $\Sigma(F)=\Sigma(P)$, and with $G_P\subset
G_F$ for the groups generated by the reflections in their facets. They
prove that $1<[G_F:G_P]<2^n$ holds exactly in three cases, the last being
``$\Sigma(F)=\Sigma(P)=\widetilde F_4$, $[G_F:G_P]=4$''. Equal Coxeter
diagrams make $P$ similar to $F$. A fundamental simplex of an index-$4$
reflection subgroup of $G_F$ is a union of $4$ fundamental simplices of
$G_F$, which are congruent. So their $\widetilde F_4$ case is a simplex
similar to the $F_4$ alcove and tiled by four alcoves: in substance,
Theorem~\ref{thm:main}. Their paper does not discuss reptiles, and we
found no reference connecting it to the reptile question. Our
contribution is that connection, the explicit tiling with a direct
proof, and its checks. The other two cases of their lemma, of types
$\widetilde C_2$ and $\widetilde G_2$, give the classical $2$-reptile
isosceles right triangle and $3$-reptile $30$-$60$-$90$ triangle.

\section{Proof of Theorem~\ref{thm:main}}\label{sec:proof}

We describe $S$ by inequalities. For $x=(x_1,x_2,x_3,x_4)$ let
\[
\ell_0(x)=1-x_1-x_2,\quad \ell_1(x)=x_2-x_3,\quad \ell_2(x)=x_3-x_4,\quad
\ell_3(x)=x_4,\quad \ell_4(x)=x_1-x_2-x_3-x_4 .
\]

\begin{lemma}\label{lem:facets}
$S=\{x\in\mathbb{R}^4:\ \ell_i(x)\ge0 \text{ for } i=0,\dots,4\}$, that is,
$S=\{x:\ x_2\ge x_3\ge x_4\ge0,\ x_1\ge x_2+x_3+x_4,\ x_1+x_2\le1\}$.
\end{lemma}

\begin{proof}
Each vertex of $S$ makes exactly one of the five functions nonzero, and
that one positive: $\ell_0(v_0)=1$, $\ell_4(v_1)=1$, $\ell_1(v_2)=\tfrac12$,
$\ell_2(v_3)=\tfrac13$ and $\ell_3(v_4)=\tfrac14$, and every other value
$\ell_i(v_j)$ is $0$. The points $v_j$ are affinely independent, since
$v_1,\dots,v_4$ are linearly independent. Let $\lambda_0,\dots,\lambda_4$
be the barycentric coordinates with respect to $v_0,\dots,v_4$. For each
vertex there is one function, say $\ell_{\sigma(j)}$ for $v_j$, that is
nonzero on it; the affine functions $\ell_{\sigma(j)}/\ell_{\sigma(j)}(v_j)$
and $\lambda_j$ agree on all five vertices, hence everywhere. So
$S=\{\lambda_j\ge0\ \forall j\}=\{\ell_i\ge0\ \forall i\}$.
\end{proof}

Let $\varphi(x)=(x_1+x_2,\ x_1-x_2,\ x_3+x_4,\ x_3-x_4)$. Its matrix
$\Phi$ satisfies $\Phi^{\mathsf T}\Phi=2I$, so $\varphi/\sqrt2$ is an
orthogonal map and $T=\varphi(S)$ is similar to $S$ with ratio $\sqrt2$;
also $\det\Phi=4$. Substituting $x_1=\frac{y_1+y_2}2$,
$x_2=\frac{y_1-y_2}2$, $x_3=\frac{y_3+y_4}2$, $x_4=\frac{y_3-y_4}2$ into
Lemma~\ref{lem:facets} gives
\begin{equation}\label{eq:T}
T=\{y\in\mathbb{R}^4:\ y_1\le1,\ y_1\ge y_2+y_3+y_4,\ y_2\ge y_3\ge y_4\ge0\}.
\end{equation}
Its vertices are $\varphi(v_0)=(0,0,0,0)$, $\varphi(v_1)=(1,1,0,0)$,
$\varphi(v_2)=(1,0,0,0)$, $\varphi(v_3)=(1,\tfrac13,\tfrac13,\tfrac13)$
and $\varphi(v_4)=(1,\tfrac12,\tfrac12,0)$. Note that $S\subset T$ by
Lemma~\ref{lem:facets} and~\eqref{eq:T}.

Define four isometries of $\mathbb{R}^4$:
\begin{align*}
g_0(x)&=x, & g_1(x)&=(1-x_2,\,1-x_1,\,x_3,\,x_4),\\
g_2(x)&=(1-x_3,\,1-x_1,\,x_2,\,x_4), & g_3(x)&=(1-x_4,\,1-x_1,\,x_2,\,x_3).
\end{align*}
Each is a signed permutation of coordinates followed by a translation.
The linear parts of $g_0,g_1,g_2,g_3$ have determinants $1,-1,1,-1$.

\begin{lemma}\label{lem:pieces}
Let $P_i=g_i(S)$. Then
\begin{align*}
P_0&=T\cap\{y:\ y_1+y_2\le1\}, &
P_1&=T\cap\{y:\ y_1+y_3\le1\le y_1+y_2\},\\
P_2&=T\cap\{y:\ y_1+y_4\le1\le y_1+y_3\}, &
P_3&=T\cap\{y:\ 1\le y_1+y_4\}.
\end{align*}
\end{lemma}

\begin{proof}
Since $g_i$ is a bijection, $y\in P_i$ exactly when $x=g_i^{-1}(y)$
satisfies the five inequalities of Lemma~\ref{lem:facets}. We write them
out in $y$ for each $i$.

For $i=0$: $y_2\ge y_3\ge y_4\ge0$, $y_1\ge y_2+y_3+y_4$ and
$y_1+y_2\le1$. The remaining inequality $y_1\le1$ of~\eqref{eq:T}
follows from $y_1\le1-y_2\le1$.

For $i=1$, put $x=(1-y_2,\,1-y_1,\,y_3,\,y_4)$. The five inequalities
become $y_1+y_3\le1$, $y_3\ge y_4$, $y_4\ge0$, $y_1\ge y_2+y_3+y_4$ and
$y_1+y_2\ge1$. The inequalities $y_1\le1$ and $y_2\ge y_3$ of~\eqref{eq:T}
follow: $y_1\le1-y_3\le1$, and $y_1+y_2\ge1\ge y_1+y_3$.

For $i=2$, put $x=(1-y_2,\,y_3,\,1-y_1,\,y_4)$. The inequalities become
$y_1+y_3\ge1$, $y_1+y_4\le1$, $y_4\ge0$, $y_1\ge y_2+y_3+y_4$ and
$y_2\ge y_3$. The inequalities $y_3\ge y_4$ and $y_1\le1$ follow from
$y_1+y_3\ge1\ge y_1+y_4$ and $y_1\le1-y_4\le1$.

For $i=3$, put $x=(1-y_2,\,y_3,\,y_4,\,1-y_1)$. The inequalities become
$y_3\ge y_4$, $y_1+y_4\ge1$, $y_1\le1$, $y_1\ge y_2+y_3+y_4$ and
$y_2\ge y_3$. The inequality $y_4\ge0$ follows from $y_4\ge1-y_1\ge0$.

In each case the inequalities obtained, together with the ones shown to
follow, are exactly those of~\eqref{eq:T} and of the stated region.
\end{proof}

\begin{proof}[Proof of Theorem~\ref{thm:main}]
On $T$ we have $y_2\ge y_3\ge y_4$, so $y_1+y_2\ge y_1+y_3\ge y_1+y_4$.
For $y\in T$, the number $1$ lies above $y_1+y_2$, between two consecutive
terms of this chain, or below $y_1+y_4$, so $y$ lies in at least one of
the four regions of Lemma~\ref{lem:pieces}. Hence
$T=P_0\cup P_1\cup P_2\cup P_3$. Within $T$, two different regions intersect only
inside one of the hyperplanes $y_1+y_j=1$ ($j=2,3,4$), so the $P_i$ have
pairwise disjoint interiors. Each $P_i=g_i(S)$ is congruent to $S$.

Apply $\varphi^{-1}$. The sets $S_i=\varphi^{-1}(P_i)$ tile
$\varphi^{-1}(T)=S$, and $S_i$ is the image of $S$ under the similarity
$\varphi^{-1}\circ g_i$ of ratio $1/\sqrt2$. These four similarities
differ from one another by isometries, so the $S_i$ are mutually
congruent. Since $\det\Phi>0$, the similarity $\varphi^{-1}\circ g_i$
preserves orientation exactly when $g_i$ does, that is, for $i=0,2$.
\end{proof}

The tiling is a chain: consecutive pieces share a facet. The vertices
of the pieces are listed in Appendix~\ref{app:vertices}.

\begin{remark}\label{rem:chiral}
The mirror images cannot be avoided by relabeling: $S$ has no
nontrivial symmetry, since its ten edge lengths admit no nonidentity
permutation of the vertices preserving them. Whether some other tiling
of $S$ by four copies uses only directly congruent ones is not settled
here.
\end{remark}

\section{Where the tiling comes from}\label{sec:origin}

Take the root system of type $F_4$ in Bourbaki's coordinates
\cite[Plate~VIII]{Bou}: its $48$ roots are $\pm e_i$, $\pm e_i\pm e_j$
($i<j$) and $\frac12(\pm e_1\pm e_2\pm e_3\pm e_4)$, with simple roots
$e_2-e_3$, $e_3-e_4$, $e_4$, $\frac12(e_1-e_2-e_3-e_4)$ and highest root
$e_1+e_2$. The fundamental alcove $\{x:\ \langle\alpha,x\rangle\ge0
\text{ for the simple roots},\ \langle e_1+e_2,x\rangle\le1\}$ is exactly
the description of $S$ in Lemma~\ref{lem:facets}. The affine Weyl group
is generated by the reflections in the hyperplanes $\langle\alpha,x\rangle
=n$ ($\alpha$ a root, $n\in\mathbb{Z}$), and its alcoves, the closures of
the regions of this arrangement, are congruent copies of $S$ tiling
$\mathbb{R}^4$. The matrix $\Phi$ is symmetric with $\Phi^2=2I$, so
$\varphi^{-1}(y)=\Phi y/2$ and $\varphi$ maps the hyperplane
$\langle\alpha,x\rangle=n$ onto $\langle\Phi\alpha,y\rangle=2n$. A direct
check of the $48$ roots shows that $\Phi$ maps each short root
($\pm e_i$ and $\frac12(\pm e_1\pm e_2\pm e_3\pm e_4)$) to a long root and
each long root to twice a short root. So the image hyperplane is
$\langle\beta,y\rangle=2n$ or $\langle\beta,y\rangle=n$ for a root $\beta$,
again a hyperplane of the arrangement. Hence $T=\varphi(S)$, the image of
a region of the arrangement, is a union of alcoves, $\det\Phi=4$ of them
by volume. These are the pieces $P_i$, and the third program in
Section~\ref{sec:verification} confirms that each is an alcove. This is
the $\widetilde F_4$ case of \cite[Lemma~10]{FT}, written in coordinates.

\begin{corollary}\label{cor:powers}
For all integers $a\ge0$ and $m\ge1$, $S$ is a $4^am^4$-reptile.
\end{corollary}

\begin{proof}
Iterating Theorem~\ref{thm:main} $a$ times tiles $S$ by $4^a$ congruent
copies scaled by $2^{-a/2}$. The dilate $mS$ is the set
$\{x:\ \langle\alpha,x\rangle\ge0 \text{ for the simple roots},\
\langle e_1+e_2,x\rangle\le m\}$, whose facets lie in hyperplanes of the
arrangement, so it is a union of alcoves, $m^4$ of them by volume. Scaling
by $1/m$ tiles $S$ by $m^4$ copies of $S/m$; apply this inside each of the
$4^a$ pieces.
\end{proof}

For $a$ odd, $4^am^4$ is a square but not a fourth power. Whether $k=9$,
$25$, $36,\dots$ ($m=3,5,6,\dots$ in the question of Kynčl and Patáková)
occur in dimension four, and whether any $k$-reptile simplex with
$k\ne m^d$ exists for $d\ge5$, remain open.

\section{Computational verification}\label{sec:verification}

The proof above is complete without a computer. The same tiling is also
checked in exact rational arithmetic by three programs, which read a
certificate file listing $S$, $\Phi$ and, for each piece, an orthogonal
matrix $M_i$, a translation $t_i$ (so that $P_i=M_iS+t_i$) and the five
vertices, together with one separating hyperplane for each pair of
pieces.
\begin{enumerate}
\item The first program checks that $S$ is nondegenerate,
  $\Phi^{\mathsf T}\Phi=2I$, each $M_i$ is orthogonal and the listed vertices
  are $M_iS+t_i$, every piece vertex has nonnegative barycentric
  coordinates in $T$, each pair of pieces lies on opposite closed sides
  of its listed hyperplane, and the four volumes add up to
  $\operatorname{vol}T=\frac1{144}=4\operatorname{vol}S$.
\item The second, written separately from the first, compares $S$ and
  $\Phi$ with the statement above, ignores the listed hyperplanes, checks
  congruence by the ten squared edge lengths as well as by $M_i$, checks
  containment by the facet inequalities of $T$, and shows that each
  pairwise intersection has dimension at most $3$ by enumerating its
  vertices exactly; then it checks the volume sum.
\item The third evaluates all $48$ roots of $F_4$ on the vertices and
  checks that $S$ and each piece lie in a closed alcove, that the pieces
  are distinct, have the volume of $S$ and lie in $T$, and that
  $\operatorname{vol}T=4\operatorname{vol}S$. As $S$ is an alcove, the
  pieces are then four distinct alcoves in $T$.
\end{enumerate}
Closed convex pieces inside $T$ with disjoint interiors and total volume
$\operatorname{vol}T$ cover $T$, so each of the first two programs alone
verifies the tiling of $T$ and hence Theorem~\ref{thm:main}. Six
corrupted certificates (a moved vertex, a map that is not a similarity,
a duplicated piece, a dropped piece, a shifted separating hyperplane, and
a non-orthogonal $M_i$) are each rejected by every program that reads the
corrupted field. The whole check takes under a second. The written proof
of Section~\ref{sec:proof} does not depend on these programs; they check
the same four pieces. The certificate, the programs and the command that runs them are
in the folder accompanying this paper, whose README describes them.

\begin{appendices}
\section{The vertices of the pieces}\label{app:vertices}

With $v_2,v_3,v_4$ as in the introduction and $w_1=(1,1,0,0)$,
$w_2=(1,\tfrac12,\tfrac12,0)$, $w_3=(1,\tfrac13,\tfrac13,\tfrac13)$, the
pieces $P_i=g_i(S)$ of $T$ are
\begin{align*}
P_0&=\operatorname{conv}\{0,\,v_1,\,v_2,\,v_3,\,v_4\}=S, &
P_1&=\operatorname{conv}\{w_1,\,v_1,\,v_2,\,v_3,\,v_4\},\\
P_2&=\operatorname{conv}\{w_1,\,v_1,\,w_2,\,v_3,\,v_4\}, &
P_3&=\operatorname{conv}\{w_1,\,v_1,\,w_2,\,w_3,\,v_4\}.
\end{align*}
Each piece arises from the previous one by replacing one vertex, and
$T=\operatorname{conv}\{0,v_1,w_1,w_2,w_3\}$. The points $v_2,v_3,v_4$
lie on the edges of $T$ from $0$: $v_2=\frac12w_1$, $v_3=\frac23w_2$,
$v_4=\frac34w_3$. The tiling of $S$ in Theorem~\ref{thm:main} is the
image of this one under $\varphi^{-1}$.

\end{appendices}

\begin{thebibliography}{9}

\bibitem{Bou}
N.~Bourbaki, \emph{Lie Groups and Lie Algebras, Chapters 4--6},
Springer, 2002, Plate VIII (system of type $F_4$).

\bibitem{FT}
A.~A. Felikson and P.~V. Tumarkin, Reflection subgroups of Euclidean
reflection groups, \emph{Sb. Math.} \textbf{196} (2005), 1349--1369;
\href{https://arxiv.org/abs/math/0402403}{arXiv:math/0402403} (February
25, 2004), Lemma~10.

\bibitem{He}
E.~Hertel, A class of self-similar $d$-polytopes, \emph{Rev. Roumaine
Math. Pures Appl.} \textbf{50} (2005), 647--655.

\bibitem{KP}
J.~Kynčl and Z.~Patáková, On the nonexistence of $k$-reptile simplices
in $\mathbb{R}^3$ and $\mathbb{R}^4$, \emph{Electron. J. Combin.}
\textbf{24}(3) (2017), \#P3.1,
\href{https://doi.org/10.37236/6113}{doi:10.37236/6113};
\href{https://arxiv.org/abs/1602.04668}{arXiv:1602.04668}.

\bibitem{MS}
J.~Matoušek and Z.~Safernová, On the nonexistence of $k$-reptile
tetrahedra, \emph{Discrete Comput. Geom.} \textbf{46} (2011), 599--609,
\href{https://doi.org/10.1007/s00454-011-9334-z}{doi:10.1007/s00454-011-9334-z};
\href{https://arxiv.org/abs/1006.1807}{arXiv:1006.1807}.

\end{thebibliography}
\end{document}
